% Auto-generated by scripts/compute_normalization_stats.py
\begin{table}[t]
\centering
\small
\begin{tabular}{@{}lrrrrrc@{}}
\toprule
\textbf{Dimension} & \textbf{Raw} & \textbf{Canon.} & \textbf{Comp.} & \textbf{Top-10} & \textbf{Top-30} & \textbf{Conv.} \\
 & \textbf{ent.} & \textbf{forms} & \textbf{ratio} & \textbf{cov.\%} & \textbf{cov.\%} & \\
\midrule
Flavor & 2,705 & 170 & 15.9$\times$ & 39.7 & 58.7 &  \\
Texture & 1,003 & 108 & 9.3$\times$ & 49.8 & 74.8 &  \\
Meal Type & 385 & 39 & 9.9$\times$ & 80.0 & 96.2 & \checkmark \\
Occasion & 815 & 56 & 14.6$\times$ & 56.6 & 86.3 & \checkmark \\
Audience & 917 & 100 & 9.2$\times$ & 63.1 & 82.1 & \checkmark \\
Cuisine Type & 251 & 20 & 12.6$\times$ & 86.6 & 100.0 & \checkmark \\
\bottomrule
\end{tabular}
\caption{Effectiveness of the agentic normalization loop. \emph{Comp.\ ratio} is the raw-to-canonical fragmentation reduction. \emph{Top-30 cov.} is the convergence criterion of the loop (target $\geq 80\%$); a checkmark in \emph{Conv.} indicates the dimension converged. Dimensions with broader expressive space (Flavor, Texture) plateau below 80\%, consistent with Zipfian long tails in opinion text~\cite{zhang2022survey}.}
\label{tab:normalization}
\end{table}
